% =====================================================================
%  3 UŽDUOTIS — DI metodų parinkimas ir pagrindimas
%  Rugsėjo 3 d. | Tikslinė apimtis: 3,5–4 psl. (RIBA 4,5)
%
%  REIKALINGA PREAMBULĖJE: \usepackage{xltabular}
%  DĖMESIO: šiame faile \section BŪTI NEGALI — jį deklaruoja
%  ataskaita.tex (rugs. 2 d. pamoka: dubliuotas \section nustūmė
%  visą turinį į kitą skyrių).
%
%  Būklė: tab:filtras užpildyta (T2). Toliau — tab:kriterijai (T3),
%  tab:matrica (T4), jautrumo analizė (T5), protokolas (T6).
%  Poskyrio 3.1 tekstas rašomas PASKUTINIS.
% =====================================================================

\subsection{Atrankos rėmas}
\label{sec:remas}

Kandidatų aibė siaurinama dviem pakopomis. Pirma pakopa yra kietieji
apribojimai. Metodas arba taikytinas šiems duomenims ir šiai diegimo
vietai, arba ne, todėl svorių ir balų čia nėra. Antra pakopa yra
svertiniai balai, taikomi tik pirmąją įveikusiems metodams.

\subsection{Kietieji apribojimai}
\label{sec:filtras}

% ---------------------------------------------------------------------
%  tab:filtras — kietųjų apribojimų filtras (T2)
% ---------------------------------------------------------------------
%  Eilučių tvarka SUTAMPA su tab:metodai (18 eilučių, tos pačios
%  paradigmų grupės) — kad skaitytojas galėtų gretinti abi lenteles.
%
%  Vartai taikomi iš eilės K1 -> K4; sprendžiantys yra pirmieji, ties
%  kuriais atsiranda x. Pažymimi VISI neįveikti vartai, ne tik pirmas.
% ---------------------------------------------------------------------

\begingroup
\footnotesize
\setlength{\tabcolsep}{3pt}
\sloppy
\begin{xltabular}{\textwidth}{@{}
  >{\raggedright\arraybackslash}p{3.5cm}
  >{\centering\arraybackslash}p{0.8cm}
  >{\centering\arraybackslash}p{0.8cm}
  >{\centering\arraybackslash}p{0.8cm}
  >{\centering\arraybackslash}p{0.8cm}
  >{\raggedright\arraybackslash}X@{}}
\caption{Kietųjų apribojimų filtras: kurie metodai apskritai taikytini
         šio darbo duomenims ir diegimo vietai}
\label{tab:filtras} \\
\toprule
\textbf{Metodas} & \textbf{K1} & \textbf{K2} & \textbf{K3} & \textbf{K4}
  & \textbf{Rezultatas} \\
\midrule
\endfirsthead
\multicolumn{6}{@{}l}{\itshape \ref{tab:filtras} lentelės tęsinys} \\
\toprule
\textbf{Metodas} & \textbf{K1} & \textbf{K2} & \textbf{K3} & \textbf{K4}
  & \textbf{Rezultatas} \\
\midrule
\endhead
\bottomrule
\endlastfoot

% --- Prižiūrimas mokymasis ---
Sprendimų medis (CART) & $\surd$ & $\surd$ & $\surd$ & $\surd$ &
  \textbf{Praeina} \\
Random Forest & $\surd$ & $\surd$ & $\surd$ & $\surd$ &
  \textbf{Praeina} \\
XGBoost & $\surd$ & $\surd$ & $\surd$ & $\surd$ &
  \textbf{Praeina} \\
LightGBM & $\surd$ & $\surd$ & $\surd$ & $\surd$ &
  \textbf{Praeina} \\
SVM (RBF) & $\surd$ & $\surd$ & $\times$ & $\surd$ &
  Atmetama (K3): $O(n^2)$--$O(n^3)$ mokymas prie 2,43~mln. eilučių \\
$k$ artimiausių kaimynų & $\surd$ & $\surd$ & $\surd$ & $\times$ &
  Atmetama (K4): inferencijos kaina auga su mokymo aibe \\
Naive Bayes & $\surd$ & $\times$ & $\surd$ & $\surd$ &
  Atmetama (K2): nepriklausomumo prielaida paneigta patikra \\
\addlinespace
% --- Neprižiūrimas mokymasis ---
Autokoderis & $\surd$ & $\surd$ & $\surd$ & $\surd$ &
  \textbf{Praeina} \\
Isolation Forest & $\surd$ & $\surd$ & $\surd$ & $\surd$ &
  \textbf{Praeina} \\
One-Class SVM & $\surd$ & $\surd$ & $\times$ & $\surd$ &
  Atmetama (K3): ta pati mokymo kaina kaip SVM \\
Klasterizavimas ($k$-vid., DBSCAN) & $\surd$ & $\surd$ & $\surd$ &
  $\surd$ & Už darbo ribų: duoda grupes, ne sprendimą
  „ataka / ne ataka“ \\
\addlinespace
% --- Pusiau prižiūrimas ---
Savimoka (pseudo-žymėjimas) & $\surd$ & $\surd$ & $\surd$ & $\surd$ &
  Už darbo ribų: rinkinys pažymėtas visas, todėl paradigmos nauda
  neišmatuojama \\
\addlinespace
% --- Gilusis mokymasis ---
Daugiasluoksnis perceptronas & $\surd$ & $\surd$ & $\surd$ & $\surd$ &
  \textbf{Praeina} \\
1D-CNN & $\surd$ & $\surd$ & $\surd$ & $\surd$ &
  \textbf{Praeina} \\
LSTM / GRU & $\times$ & $\surd$ & $\surd$ & $\surd$ &
  Atmetama (K1): sekos duomenyse nėra, nes kiekviena eilutė jau yra
  lango agregatas \\
Transformer & $\times$ & $\surd$ & $\times$ & $\times$ &
  Atmetama (K1): ta pati sekos problema; papildomai netelpa į mokymo
  ir delsos biudžetus \\
Grafų neuroninis tinklas & $\times$ & $\surd$ & $\times$ & $\times$ &
  Atmetama (K1): nėra mazgų identifikatorių nei briaunų, grafo
  sukonstruoti neįmanoma \\
\addlinespace
% --- Paskirstytas mokymasis ---
Federuotas mokymasis (FedAvg) & $\times$ & $\surd$ & $\surd$ & $\surd$ &
  Atmetama (K1): nėra įrenginio identifikatoriaus, realistiško ne-IID
  padalijimo į klientus sudaryti neįmanoma \\
\end{xltabular}

\vspace{-0.5em}
{\footnotesize
Vartai.
\textbf{K1}: duomenų struktūra leidžia metodą taikyti;
\textbf{K2}: modelio prielaidos nepažeistos duomenimis;
\textbf{K3}: mokymas telpa į \SI{30}{\minute} be grafinio
spartintuvo prie 2,43~mln. eilučių ir 36 požymių;
\textbf{K4}: inferencija telpa į \SIrange{20}{50}{\milli\second}
kraštinio šliuzo biudžetą \cite{sallam2026gap}.
$\surd$ žymi įveiktus vartus, $\times$ neįveiktus.
\par}
\endgroup

\clearpage
\subsection{Vertinimo kriterijai ir jų svoriai}
\label{sec:kriterijai}

Kiekvienas antrosios pakopos kriterijus nurodo
\ref{tab:reikalavimai} lentelės eilutę, iš kurios kyla. Kartu su
kriterijumi apibrėžiama ir balų skalė, todėl balai turi vienodą reikšmę
visiems metodams.

% ---------------------------------------------------------------------
%  tab:kriterijai — kriterijai, svoriai, kilmė ir balų skalė (T3)
% ---------------------------------------------------------------------
%  PAKEITĖ atšauktą tab:reikalavimai perkėlimą iš 1 skyriaus:
%  lentelė lieka savo vietoje, čia į ją tik nurodoma.
%
%  Stulpelis „Iš ko kyla“ yra visa šios lentelės vertė — be jo
%  svoriai būtų nuomonė. Balų skalė apibrėžiama ČIA, prieš balus
%  (tab:matrica), o ne po jų.
% ---------------------------------------------------------------------

\begin{table}[H]
\centering
\caption{Metodų vertinimo kriterijai, jų svoriai ir kilmė}
\label{tab:kriterijai}
\footnotesize
\setlength{\tabcolsep}{4pt}
\begin{tabularx}{\textwidth}{@{}
  >{\raggedright\arraybackslash}p{2.9cm}
  >{\centering\arraybackslash}p{0.9cm}
  >{\raggedright\arraybackslash}X
  >{\raggedright\arraybackslash}X@{}}
\toprule
\textbf{Kriterijus} & \textbf{Svoris} & \textbf{Iš ko kyla}
  & \textbf{Balų skalė} \\
\midrule

Aptikimo kokybė (macro-F1 potencialas)
  & 30~\%
  & \ref{tab:reikalavimai} lentelės ketvirta eilutė: pažeidžiamumai
    lieka neištaisyti visą įrenginio tarnavimo laiką, todėl aptikimas
    yra kompensacinė priemonė ir vienintelė likusi apsauga
  & 5: ant CICIoT2023 skelbiama macro-F1 $\geq 0{,}89$
    \cite{almahaqeri2026gradient}; 1: žemiau $0{,}70$ \\
\addlinespace

Klaidingi teigiami ir atsparumas disbalansui
  & 30~\%
  & \ref{tab:reikalavimai} lentelės penkta eilutė ir \ref{sec:atakos}
    skyriuje priimta 1~\% klaidingų teigiamų riba
  & 5: disbalansas sprendžiamas klasių svoriais, nedubliuojant
    duomenų; 1: mokymo metu disbalansas blogėja \\
\addlinespace

Resursų poreikis (inferencijos delsa, modelio dydis)
  & 25~\%
  & \ref{tab:reikalavimai} lentelės pirma eilutė, kurioje resursai
    įvardyti kaip atrankos kriterijus
  & 5: inferencija mikrosekundžių eilės, modelis $<$~10~MB;
    1: delsa artėja prie \SI{20}{\milli\second} biudžeto \\
\addlinespace

Interpretuojamumas
  & 15~\%
  & Ta pati penkta eilutė. Ribotas analitikų pajėgumas reiškia, kad
    signalą reikia mokėti pagrįsti
  & 5: savaiminis; 3: per XAI su papildoma kaina
    \cite{mohale2025xai}; 1: juodoji dėžė \\
\bottomrule
\end{tabularx}

\vspace{2pt}
\raggedright\footnotesize
Likusios dvi \ref{tab:reikalavimai} lentelės eilutės kriterijaus
neduoda, nes jas vienodai tenkina visi po filtro likę metodai.
\normalsize
\end{table}

\clearpage
\subsection{Sprendimų matrica}
\label{sec:matrica}

Kiekvienam po filtro likusiam metodui skiriamas balas pagal kiekvieną
kriterijų, o svertinė suma juos surikiuoja. Prižiūrimi ir neprižiūrimi
metodai rikiuojami atskirai.

% ---------------------------------------------------------------------
%  tab:matrica — svertiniai balai aštuoniems po filtro likusiems (T4)
% ---------------------------------------------------------------------
%  LENTELĖ GENERUOJAMA, ne rašoma ranka:
%      rezultatai/darbiniai/sprendimu_matrica.csv
%      -> python -m src.eksperimentai.matrica
%      -> ataskaita/lenteles/matrica.tex
%  Perskaičiavus svorius (T5) skaičiai atsinaujina patys.
%
%  Balų skalė apibrėžta tab:kriterijai — PRIEŠ balus, ne po jų.
% ---------------------------------------------------------------------

\begin{table}[H]
\centering
\caption{Sprendimų matrica: po filtro likę metodai pagal
         \ref{tab:kriterijai} lentelės kriterijus ir svorius}
\label{tab:matrica}
\lentele{matrica.tex}

\vspace{2pt}
\raggedright\footnotesize
Balai 1--5 pagal \ref{tab:kriterijai} lentelės skalę.
Abiejų blokų sumos tarpusavyje nepalyginamos, nes prižiūrimi metodai
vertinami aštuonių kategorijų klasifikavimo, o neprižiūrimi dvejetainėje
anomalijų aptikimo formuluotėje.
Neprižiūrimų metodų stulpelis „Disbalansas“ žymi, kad klasių disbalanso
problema jiems nekyla (mokoma tik iš gerybinio srauto), o ne kad
ji išspręsta.
\normalsize
\end{table}

\subsection{Jautrumo analizė}
\label{sec:jautrumas}

% ---------------------------------------------------------------------
%  tab:jautrumas — svorių jautrumas (T5), GENERUOJAMA
%      python -m src.eksperimentai.jautrumas
%  Procedūra užrašyta plane PRIEŠ paleidžiant, kartu su abiem galimais
%  rezultatais — kad nekiltų pagundos derinti svorius prie atsakymo.
% ---------------------------------------------------------------------

Kiekvienam kriterijui atskirai pridėta ir atimta po 10 procentinių
punktų, likusius svorius perskirstant proporcingai. Visais aštuoniais
atvejais pirmi trys prižiūrimi metodai išliko tie patys, tad papildomai
nustatyta, kokio svorio reikėtų kiekvienai porai apversti.

\begin{table}[H]
\centering
\caption{Jautrumo analizė: kada rikiuotė apsiverstų}
\label{tab:jautrumas}
\lentele{jautrumas.tex}

\vspace{2pt}
\raggedright\footnotesize
Stulpelis „Skirtumas“ rodo svertinių sumų skirtumą prie bazinių svorių.
Dominavimas reiškia, kad pirmasis metodas nėra blogesnis nė pagal
vieną kriterijų, tad tokios poros rezultato joks svorių derinys
nekeičia.
\normalsize
\end{table}

Ketverto metodų tarpusavio rikiuotė nuo svorių nepriklauso, nes
XGBoost prieš Random Forest ir Random Forest prieš daugiasluoksnį
perceptroną yra dominavimo atvejai. Apsiverstų tik poros su sprendimų
medžiu, kuris į ketvertą nepateko; XGBoost jis aplenktų tik tada, jei
interpretuojamumui būtų skiriama beveik tris kartus daugiau svorio nei
dabar.

\subsection{Eksperimento protokolas}
\label{sec:protokolas}

% ---------------------------------------------------------------------
%  3.6 — protokolas UŽRAKINAMAS prieš eksperimentus (T6).
%  Pilnas 24 punktų sąrašas — claude/uzduotis_03_planas.md 5 sk.;
%  čia tik tai, kas keičia rezultatų interpretaciją.
%  TEKSTU, ne ketvirta lentele — apimties biudžetas.
% ---------------------------------------------------------------------

Eksperimento sąlygos fiksuojamos prieš pirmąjį mokymą.

Duomenys. Naudojama stratifikuota imtis, ne daugiau kaip
\num{100000} eilučių klasei; retos klasės imamos visos. Pirma šalinamos
nutrūkusios eilutės, tada begalinės \texttt{Rate} reikšmės. Iš 39
požymių šalinami trys tiksliai išvestiniai, lieka 36.

Dublikatai. Tikslūs dublikatai šalinami pagal visą eilutę ir prieš
imties sudarymą, nes atsitiktinai skaidant tos pačios eilutės kopijos
patektų ir į mokymo, ir į testavimo aibę.

Vienodi požymių vektoriai, pažymėti skirtingomis etiketėmis, paliekami
kaip tikras duomenų dviprasmiškumas. Iš jų kyla teorinė tikslumo riba,
žemesnė už 100~\%, ir už ją aukštesnis rezultatas reikštų nutekėjimą.
Riba galioja šiam 39 požymių leidimui ir dublikatų neturinčiai imčiai.

Skaidymas. Stratifikuotas atsitiktinis \num{70}/\num{15}/\num{15} su
fiksuotu pradiniu dydžiu; indeksai išsaugomi ir įkeliami. Chronologinis
skaidymas negalimas, nes rinkinyje nėra laiko žymos, todėl rezultatai
nieko nesako apie elgseną laikui bėgant. Normalizavimas ir klasių
balansavimas taikomi tik mokymo aibei.

Formuluotė. Pagrindinis detalumas yra aštuonios atakų kategorijos;
dvejetainė ir trisdešimt keturių klasių formuluotės atliekamos tik
geriausiam modeliui.

Vertinimas. Pagrindinė metrika yra macro-F1; šalia jos per-klasę
tikslumas ir atkūrimas, PR-AUC, MCC ir sumaišymo matrica. Bendras
tikslumas pateikiamas tik palyginimui su literatūra. Inferencijos delsa
matuojama gryna, atskirai nuo srauto lango sukaupimo laiko.

Disbalansas sprendžiamas klasių svoriais, nes jie nedubliuoja duomenų;
SMOTE \cite{chawla2002smote} tikrinamas kaip abliacija vienam modeliui.
Kiekvienas modelis mokomas su trimis pradiniais dydžiais, pateikiamas
vidurkis ir standartinis nuokrypis. Trijų paleidimų statistiniam testui
nepakanka \cite{dietterich1998tests}, todėl skirtumas laikomas
reikšmingu tik viršijęs paleidimų sklaidą. Hiperparametrai derinami
atsitiktine paieška validacijos aibėje, o testavimo aibė neliečiama iki
vertinimo etapo.

Nematytų atakų patikra atliekama trimis scenarijais. Iš mokymo aibės
pašalinama po vieną klasę iš daugiaklasės kategorijos
(\texttt{DDOS-SLOWLORIS}, \texttt{RECON-PORTSCAN}) ir viena ištisa
kategorija (\texttt{DICTIONARYBRUTEFORCE}, vienintelė savo kategorijos
klasė). Pirmieji du tikrina, ar atpažįstamas naujas žinomos šeimos
variantas; trečiasis yra autokoderio hipotezės patikra, nes prižiūrimas
modelis tokios kategorijos etiketės neturi.

\subsection{Galutinė aibė}
\label{sec:parinkimas_isvados}

Rikiuotė (\ref{tab:matrica} lentelė) galutinės aibės neduoda, nes imami
ne keturi geriausi balai.
Aibė sudaroma pagal rikiuotės viršūnę, paradigmų padengimą ir
sudėtingumo gradientą.

Pasirenkamas XGBoost, kuris vienintelis turi tiesioginį atitikmenį
literatūroje; Random Forest kaip ansamblio pakopa tarp medžio ir
stiprinimo, nustatanti atskaitos lygį; daugiasluoksnis perceptronas,
atsakantis, ar sudėtingumas tokiai įvesčiai apsimoka; ir autokoderis,
įtraukiamas dėl funkcinio reikalavimo aptikti atakas be jų pavyzdžių.

LightGBM balais lygus XGBoost, tad du beveik tapatūs stiprinimo metodai
palyginime nieko neprideda. Sprendimų medis matricoje lenkia Random
Forest, o Isolation Forest lenkia autokoderį, tačiau ketverte jų vietos
atitenka ansamblio pakopai ir funkciniam reikalavimui.
